\begin{example}
    Sea $a \neq 0$ y sea $\sigma > 0$ tal que $\sigma \notin \mathbb{Z}$.
    Entonces, la sucesión $(an^\sigma)_{n\geq1}$ es equidistribuida módulo $1$.
\end{example}

\begin{resolve}
    Sea $g(x) = a x^\sigma$ y sea $k = \floor{\sigma}+1$. La derivada de orden $k$ de $g$ es 
    \[
        g^{(k)}(x) = a \sigma (\sigma - 1) \cdots (\sigma - k + 1) x^{\sigma - k}.
    \]
    Como $\sigma-k < 0$,
    \[
        \lim_{x \to \infty} g^{(k)}(x) = 0.
    \]

    Para el otro límite, observemos que $\sigma-k+1 = \sigma - \floor{\sigma} > 0$, y por tanto
    \[
        \lim_{x \to \infty} x \left|g^{(k)}(x)\right| = 
        \lim_{x \to \infty} \left|a \sigma (\sigma - 1) \cdots (\sigma - k + 1)x^{\sigma-k+1} \right| = \infty.
    \]
    
    Para verificar la monotonía de $g^{(k)}(x)$, calculamos su derivada,
    \[
        g^{(k+1)}(x) = a \sigma (\sigma - 1) \cdots (\sigma - k) x^{\sigma - k - 1}
    \]
    Como $x^{\sigma-k-1}$ es positivo para todo $x\geq1$, entonces $g^{(k+1)}(x)$ 
    es de signo constante para todo $x \geq 1$. En consecuencia, $g^{(k)}(x)$ es 
    monótona a partir de $x_0 = 1$.
    
    Por el Teorema \ref{teo:fejer_cap4}, la sucesión $(an^\sigma)$ es equidistribuida módulo $1$.
\end{resolve}